\section{Auxiliary Nesting and Scope}

An auxiliary may take the following auxiliary phrase as its complement. Each earlier auxiliary applies to the entire phrase after it. Read a chain from the inside outward. Brackets here show grouping; they add no spoken or written Tovian morphemes.

Within one shared-subject chain, every auxiliary carries the same person prefix. Each layer may contribute its own mood and aspect. Within a single auxiliary, mood takes scope over aspect. Case and class prefixes remain on noun phrases in their existing order.

\subsection*{Scope contrasts}

\begin{itemize}
\item \textit{apa muril}: stops seeing.
\item \textit{ashopa muril}: can [stop seeing].
\item \textit{apa [ashofa muril]}: stops [being able to see].
\item \textit{ashoya muril}: can [begin seeing].
\item \textit{aya [ashofa muril]}: begins [being able to see].
\item \textit{atlopa for}: should [stop writing].
\item \textit{apa [atlofa for]}: ceases [to be expected to write].
\item \textit{ipa [ishofa muril]}: I/we cease [being able to see].
\end{itemize}

Agreement repeats because the grammatical subject is shared. The moods and aspects need not match.

\subsection*{Can no longer be ignored}

The intended scope is cessation over possibility over a passive predicate. With \texttt{[IGNORE]} as an English placeholder for a lexical root still to be chosen:

\begin{quote}
\texttt{apa [ashopafa [IGNORE]]}
\end{quote}

\begin{itemize}
\item \textit{a-p-a}: third person + cessative + present.
\item \textit{a-sho-pa-f-a}: third person + potential + passive + simple + present.
\item \texttt{[IGNORE]}: the unchanged lexical predicate.
\end{itemize}

The truths remain the grammatical subject throughout. The inner phrase means ``can be ignored''; the outer cessative ends that possibility. ``Stop being ignored'' describes the end of actual ignoring, while ``can stop being ignored'' expresses the possibility of stopping. The chain uses one cessative; the bound \textit{pa-} inside \textit{ashopafa} marks passive voice.

``No longer'' presupposes an earlier possibility without implying that the impossibility lasts forever. This construction adds neither a new particle nor an untensed auxiliary form.

\subsection*{Chains and embedded clauses}

``I cease [being able to see]'' has one shared subject. ``I know [you left yesterday]'' contains a clause with its own subject and time. Matching person is required within one shared-subject chain, not across an entire sentence. The temporal frame \textit{tu-le-ngothr: P; Q} specifies one construction with a finite event description P (see Temporal Clauses). General complement clauses such as ``I know [you left]'' still need their own rule.

\subsection*{Temporal dependence: open question}

The examples here use present tense throughout. Matching person does not require Tovian to copy tense endings across every layer. A general rule for inner temporal reference and mixed-tense chains remains to be defined. These examples establish person agreement and scope without settling temporal dependence.
